\ifdefined\iscombined
	\lecturetitle{Introduction}
\else
\documentclass{article}
\usepackage{xparse}
\usepackage{changepage}
\usepackage{listings}
\usepackage{amsthm}
\usepackage{comment}
\usepackage{amsmath}
\usepackage{braket}
\usepackage{amsfonts}
\usepackage{sectsty}
\usepackage{hyperref}
\usepackage{fancyhdr}
\usepackage[top=1in,bottom=1in,left=1in,right=1in]{geometry}
\usepackage{float}
\usepackage{graphicx}
\usepackage{mathtools}

\date{
	\vspace{-.75em}
	{\large \today}
}

\title{
	\vspace*{-1.5em}
	{\Huge Lecture 1} \\
	\vspace{-1em}
	\rule{\textwidth/4}{.01em} \\
	\vspace{-.25em}
	{\Large Introduction}
	\vspace{-.75em}
}

\author{
	{\large Matthew Farah}
}

\hypersetup{
	colorlinks=true,
	linkcolor=blue,
	filecolor=magenta,
	urlcolor=blue,
	citecolor=blue,
	anchorcolor=blue
}

\fancyhead{}
\fancyhead[L]{\today}
\fancyhead[R]{Lecture 1 - Introduction}

\setlength{\parindent}{0cm}

\newcounter{example}
\renewcommand{\theexample}{\arabic{example}}

\newcounter{subexample}
\renewcommand{\thesubexample}{\alph{subexample}}

\newenvironment{example}[1][]{
	\refstepcounter{example}
	\setcounter{subexample}{0}
	\begin{adjustwidth}{1.5em}{}
	\noindent\textbf{\ifx&#1&Example \theexample:\else Example \theexample \ (#1):\fi}
	\ignorespaces
}{
	\vspace{.5em}
	\end{adjustwidth}
}

\newenvironment{solution}{
	\vspace{.5em}
	\par
	\begin{adjustwidth}{1.5em}{}
	\textbf{{Solution:}}
	\par
	\vspace{.5em}
}{
	\vspace{1em}
	\end{adjustwidth}
}

\newenvironment{subexample}{
	\setcounter{step}{0}
	\refstepcounter{subexample}
	\vspace{1em}\par\noindent
	\begin{adjustwidth}{1.5em}{}
	\noindent\textbf{(\thesubexample)}
	\ignorespaces
}{
	\par
	\vspace{1em}
	\end{adjustwidth}
}

\renewenvironment{proof}{
	\vspace{.5em}\par\noindent
	\begin{adjustwidth}{1.5em}{}
	\emph{Proof:}\par\vspace{1em}
	\setlength{\parindent}{0cm}
}{
	\end{adjustwidth}
}

\newtheorem{theorem}{Theorem}

\renewenvironment{theorem}[1][]{
	\refstepcounter{theorem}
	\vspace{1em}\par\noindent
	\begin{adjustwidth}{1.5em}{}
	\noindent\textbf{\ifx&#1&Theorem \thetheorem:\else Theorem \thetheorem \ (#1):\fi}
	\ignorespaces
}{
	\vspace{.5em}
	\end{adjustwidth}
}

\newtheorem{lemma}{Lemma}

\renewenvironment{lemma}[1][]{
	\refstepcounter{lemma}
	\vspace{1em}\par\noindent
	\begin{adjustwidth}{1.5em}{}
	\noindent\textbf{\ifx&#1&Lemma \thelemma:\else Lemma \thelemma \ (#1):\fi}
	\ignorespaces
}{
	\vspace{.5em}
	\end{adjustwidth}
}

\newtheorem{definition}{Definition}

\renewenvironment{definition}[1][]{
	\refstepcounter{definition}
	\vspace{1em}\par\noindent
	\begin{adjustwidth}{1.5em}{}
	\noindent\textbf{\ifx&#1&Definition \thedefinition:\else Definition \thedefinition \ (#1):\fi}
	\ignorespaces
}{
	\vspace{.5em}
	\end{adjustwidth}
}

\newtheorem{corollary}{Corollary}

\renewenvironment{corollary}[1][]{
	\refstepcounter{corollary}
	\vspace{1em}\par\noindent
	\begin{adjustwidth}{1.5em}{}
	\noindent\textbf{\ifx&#1&Corollary \thecorollary:\else Corollary \thecorollary \ (#1):\fi}
	\ignorespaces
}{
	\vspace{.5em}
	\end{adjustwidth}
}

\newtheorem{proposition}{Proposition}

\renewenvironment{proposition}[1][]{
	\refstepcounter{proposition}
	\vspace{1em}\par\noindent
	\begin{adjustwidth}{1.5em}{}
	\noindent\textbf{\ifx&#1&Proposition \theproposition:\else Proposition \theproposition \ (#1):\fi}
	\ignorespaces
}{
	\vspace{.5em}
	\end{adjustwidth}
}

\newtheorem{remark}{Remark}

\renewenvironment{remark}[1][]{
	\refstepcounter{remark}
	\vspace{1em}\par\noindent
	\begin{adjustwidth}{1.5em}{}
	\noindent\textbf{\ifx&#1&Remark \theremark:\else Remark \theremark \ (#1):\fi}
	\ignorespaces
}{
	\vspace{.5em}
	\end{adjustwidth}
}

\makeatother
\makeatletter

\newcounter{step}
\renewcommand{\thestep}{\Roman{step}}

\newenvironment{step}{
	\refstepcounter{step}
	\vspace{1em}\par\noindent
	\ignorespaces\noindent\textbf{(\thestep)}
	\ignorespaces
}{
	\par
}

\sectionfont{\fontsize{12}{12}\bfseries}
\subsectionfont{\fontsize{10}{12}\normalfont}
\subsubsectionfont{\fontsize{10}{12}\normalfont}

\begin{document}

\maketitle
\pagestyle{fancy}

\fi

\begin{itemize}
	\item Notes about the course:
	\begin{itemize}
		\item 4 problems (5\% each).
		\item Weekly quizzes (10\% total).
		\item 2 midterms (15\% each).
		\item Exam (40\% total).
	\end{itemize}
	\item Algebra is concerned with studying the implications that certain collections of axioms have on solving problems.
\end{itemize}

\begin{definition}
	Let $G$ be a set. A binary operator $*$ is a function $*: G \times G \to G$.
\end{definition}

\begin{definition}
	Let $G$ be a set. $G$ is a group under a binary operator $*$ if $(G, *)$ satisfies all 3 of the following criteria:

	\begin{enumerate}
		\item Associativity:
		\begin{itemize}
			\item $\forall a, b, c \in G; \ (a * b) * c = a * (b * c)$.
		\end{itemize}
		\item Existence of an identity element:
		\begin{itemize}
			\item $\exists e \in G; \ \forall a \in G; \ a * e = e * a = a$.
		\end{itemize}
		\item Existence of inverses:
		\begin{itemize}
			\item $\forall a \in G; \ \exists b \in G; \ a * b = b * a = e$.
		\end{itemize}
	\end{enumerate}
\end{definition}

\begin{definition}
	A group $(G, *)$ is called abelian if, in addition to satisfying the group axioms, $*$ satisfies the commutativity of $*$ on $G$:

	\begin{gather*}
		\forall a, b \in G; \ a * b = b * a
	\end{gather*}
\end{definition}

\begin{example}
	Which of the following are groups?
\end{example}

\begin{subexample}
	$\mathbb{Z}$ under addition?
\end{subexample}

\begin{solution}
	Yes!
\end{solution}

\begin{subexample}
	$\mathbb{N}$ under addition?
\end{subexample}

\begin{solution}
	No! There is no inverse for any element except $0$ if you let $0 \in \mathbb{N}$.
\end{solution}

\begin{subexample}
	$\mathbb{Z}$ under multiplication?
\end{subexample}

\begin{solution}
	No! Many elements, such as $2$, do not have a multiplicative inverse in $\mathbb{Z}$.
\end{solution}

\begin{subexample}
	$\mathbb{Z}$ under subtraction?
\end{subexample}

\begin{solution}
	No! Subtraction isn't associative. For example:

	\begin{align*}
		(1 - 2) - 3 = -4 \ne 2 = 1 - (2 - 3)
	\end{align*}
\end{solution}

\begin{example}
	Do any of the following sets under addition constitute a group?

	\begin{align*}
		\{1, 2, 3\}, \quad \{0, 1, 2, 3\}, \quad \{-3, -2, -1, 0, 1, 2, 3\}
	\end{align*}
\end{example}

\begin{solution}
	No, since none of them are closed under addition. For example, $1 + 3 = 4$ does not belong to any of the sets.
\end{solution}

\ifdefined\iscombined
	\newpage
\else
	\end{document}
\fi
